
%% ============================================================
%%  PAPER 6 — RQ6
%%  Kinematic-RCS Feature Fusion for Maritime Vessel
%%  Classification: Temporal Behavioural Signatures Break the
%%  RCS-Only Accuracy Ceiling
%% ============================================================
\documentclass[journal]{IEEEtran}
\usepackage{amsmath,amssymb,amsfonts}
\usepackage{booktabs,array,multirow,tabularx}
\usepackage{graphicx,float,caption}
\usepackage{xcolor,hyperref,cite}
\usepackage{algorithm,algpseudocode}
\hypersetup{colorlinks=true,linkcolor=blue,citecolor=blue,urlcolor=blue}

\newcommand{\vect}[1]{\boldsymbol{#1}}
\newcommand{\ket}[1]{|#1\rangle}
\newcommand{\expect}[1]{\langle #1 \rangle}
\newcommand{\pp}{\,\text{pp}}

\begin{document}

\title{
  \textbf{Kinematic--RCS Feature Fusion for Maritime\\
  Vessel Classification: Temporal Behavioural Signatures\\
  Break the RCS-Only Accuracy Ceiling}\\[0.4em]
  \large Research Question 6: Does fusing temporally-aggregated radar
  kinematic signatures with range-invariant RCS features surpass the
  classification ceiling of either feature family alone?
}
\author{Vijay~Kumar~V%
\IEEEcompsocitemizethanks{
\IEEEcompsocthanksitem V.~Kumar is an Independent Researcher and Innovator,
Bangalore, India. E-mail: vinayitsv14@gmail.com.%
\IEEEcompsocthanksitem This work was self-funded by the author as an
independent research contribution with no institutional or corporate
affiliation.%
}}

\IEEEtitleabstractindextext{%
\begin{abstract}
Papers~1--5 of this series established that range-invariant radar cross
section (RCS) features achieve 84.0--84.8\% four-class vessel
classification accuracy for Gradient Boosting Trees and 77.7--78.3\%
for Hybrid Quantum-Classical Neural Networks (HQNN). The residual error
in the RCS-only regime is dominated by the Cargo/Tanker pair, which
share overlapping hull geometry despite representing distinct vessel
archetypes. This paper introduces a \emph{kinematic augmentation} that
adds four temporally-aggregated behavioural descriptors --- instantaneous
speed-over-ground, 15-minute mean speed, 15-minute speed standard
deviation, and 15-minute course-over-ground standard deviation --- to
the five RCS features from Paper~5, forming a nine-dimensional
kinematic-RCS fused descriptor. On a 300 samples-per-class balanced
corpus of 91,205 cleaned coastal X-band radar detections (Fishing, Tug,
Cargo, Tanker), Gradient Boosting Trees achieve \textbf{94.4\% accuracy
and 94.4\% macro F1} --- the highest result in this series and a
10.4 percentage-point improvement over the RCS-only baseline.
XGBoost achieves 93.2\% (+8.4\,pp). HQNN-8q achieves 84.2\%
(+5.9\,pp), now equalling the prior classical ceiling, while HQNN-5q
achieves 81.7\% (+4.0\,pp). The fusion works because RCS features
encode \emph{what the vessel looks like} (geometry, reflectivity) while
kinematic features encode \emph{how the vessel moves} (transit speed
and course stability) --- two physically orthogonal information
channels that are conditionally nearly independent given vessel class.
Cargo and Tanker vessels, while geometrically similar at radar range
resolution, differ in their kinematic regimes: tankers maintain
exceptionally steady courses on deep-sea lanes, while cargo vessels
exhibit greater course variance during port approaches and coastal
transits. Course-over-ground standard deviation ($\sigma_\psi$) is
identified as the dominant kinematic discriminator via EDA
(Cohen's~$d = 0.357$) and confirmed in the fused model by feature
ablation. Classical ensemble methods scale substantially better than
HQNN architectures to the expanded nine-dimensional feature space,
suggesting that the quantum advantage in this problem domain lies in
calibration and uncertainty quantification rather than accuracy
maximisation when training-set size is constrained.
\end{abstract}
\begin{IEEEkeywords}
maritime vessel classification; radar; kinematic features; RCS features;
feature fusion; temporal aggregation; gradient boosting; XGBoost; quantum
machine learning; HQNN; PennyLane; behavioural signatures; Cargo-Tanker
separation; course-over-ground variance; coastal surveillance
\end{IEEEkeywords}}

\maketitle
\IEEEdisplaynontitleabstractindextext
\IEEEpeerreviewmaketitle

%% ────────────────────────────────────────────────────────────
\section{Introduction}
\label{sec:intro}

\IEEEPARstart{M}{aritime} vessel classification from terrestrial radar
faces a fundamental challenge: the radar return from a vessel encodes
both the electromagnetic interaction with the hull (size, shape,
material) and the kinematic state of the vessel (speed, course,
manoeuvrability). These two channels of information are measured
through different physical pathways and, crucially, are largely
independent given the vessel's type. A radar echo at a single instant
captures hull geometry; the sequence of echoes over a vessel's transit
captures its behavioural pattern. Using only one channel leaves
discriminative information on the table.

Papers~1--4 of this series established a rigorous baseline using five
raw geometric features (azimuth, peak amplitude, total amplitude,
down-range extent, cross-range extent in metres) and demonstrated that
Gradient Boosting Trees achieve 90.0\% accuracy on five vessel types
when trained on 300 balanced samples per class~\cite{paper1features}.
Paper~5 introduced range-invariant RCS engineering --- transforming raw
amplitudes via the monostatic radar range equation --- and achieved
84.0--84.8\% on four classes with five physically grounded features,
while demonstrating that Hybrid Quantum-Classical Neural Networks
(HQNN) achieve substantially better probability calibration than
classical ensemble methods~\cite{paper5rcs}.

Despite these advances, the \emph{Cargo-Tanker boundary} has remained
the dominant classification bottleneck. Both vessel types operate in
constrained deep-sea lanes, both have large hulls with broad radar
footprints, and their length-to-beam ratios overlap substantially.
Paper~4 identified this as the hardest pair in the four-class
problem~\cite{paper4ct}, and Paper~5 confirmed that the RCS-only
confusion rate remains around 25--30\% in either direction.

This paper tests whether the kinematic channel --- available in any
coherent radar that tracks detections over time --- provides the
complementary information needed to break this ceiling. Specifically,
we augment the five RCS features with four 15-minute rolling statistics
of speed and course behaviour, forming a nine-dimensional
\emph{kinematic-RCS} descriptor, and evaluate the same four classifier
architectures as Papers~2--5 (GBT, XGBoost, HQNN-5q, HQNN-8q) on the
same balanced experimental protocol.

\medskip
\noindent\textbf{Principal contributions of this paper:}
\begin{enumerate}
\item A kinematic-RCS fused descriptor combining two physically
      orthogonal information channels from standard coastal radar
      outputs, requiring no additional hardware or sensing modality.
\item Experimental demonstration that GBT achieves 94.4\% four-class
      accuracy on the fused descriptor --- a 10.4\,pp improvement over
      the RCS-only result of Paper~5 and the highest accuracy reported
      in this series.
\item Evidence that course-over-ground standard deviation
      ($\sigma_\psi$) is the dominant kinematic discriminator, with a
      Cohen's~$d = 0.357$ Cargo-versus-Tanker effect size confirmed by
      both EDA and feature ablation.
\item Characterisation of asymmetric scaling: classical ensemble
      methods gain 8--10\,pp from kinematic augmentation while HQNN
      architectures gain 4--6\,pp, establishing that the quantum
      advantage in this domain resides in calibration rather than in
      accuracy when training data is limited to 300 samples per class.
\item Analysis of the residual error structure after fusion: the
      dominant confusion shifts from Cargo-Tanker (geometric overlap)
      to Fishing-Tug (behavioural overlap), indicating that the fusion
      resolves the RCS-only bottleneck but reveals the next limiting
      physical ambiguity.
\end{enumerate}

The remainder of this paper is organised as follows.
Section~\ref{sec:background} summarises the relevant background.
Section~\ref{sec:features} derives the kinematic features and justifies
their physical discriminability.
Section~\ref{sec:data} describes the dataset and experimental setup.
Section~\ref{sec:results} presents quantitative results.
Section~\ref{sec:discussion} discusses the mechanisms and implications.
Section~\ref{sec:conclusion} concludes.

%% ────────────────────────────────────────────────────────────
\section{Background}
\label{sec:background}

\subsection{RCS Features from Paper~5}
\label{sec:rcs_recap}

Paper~5 derived five range-invariant features from the monostatic radar
range equation~\cite{skolnik2001}. Received power $P_r$ follows:
\begin{equation}
  P_r = \frac{P_t G^2 \lambda^2 \sigma}{(4\pi)^3 R^4}
  \label{eq:range_eq}
\end{equation}
where $P_t$ is transmit power, $G$ antenna gain, $\lambda$ wavelength,
$\sigma$ true RCS, and $R$ slant range. Taking logarithms and isolating
$\sigma$:
\begin{equation}
  \log\hat{\sigma}_p = \log A_p + 4\log R + C_{\text{sys}}
  \label{eq:rcs_peak}
\end{equation}
where $A_p$ is the per-detection peak amplitude in system units and
$C_{\text{sys}}$ is a system constant absorbed into the feature (since
classifiers are trained on the same radar). Similarly for total
amplitude:
\begin{equation}
  \log\hat{\sigma}_t = \log A_t + 4\log R + C_{\text{sys}}
  \label{eq:rcs_total}
\end{equation}
Three additional features encoding shape, concentration, and size:
\begin{align}
  \rho   &= A_p / A_t  &\text{(RCS concentration ratio)} \label{eq:conc}\\
  \xi    &= d_R / d_A  &\text{(aspect ratio, down/cross-range)} \label{eq:ar}\\
  \log A &= \log(\pi d_R d_A / 4) &\text{(log footprint area)} \label{eq:fp}
\end{align}
where $d_R$ is the down-range gate extent and $d_A$ the azimuth extent
in metres. Paper~5 showed that GBT on these five features achieves
84.0\%, with feature importance dominated by aspect ratio ($\sim$25\%)
and log footprint area ($\sim$23\%).

\subsection{Temporal Track Features in Radar Systems}

Modern pulse-Doppler maritime surveillance radars routinely maintain
vessel tracks by associating detections across successive scans.
The track processor typically estimates and reports the vessel's
\emph{speed over ground} (SoG, derived from successive Doppler
measurements or position differencing) and \emph{course over ground}
(CoG, derived from successive azimuth and range estimates). These
quantities are reported per detection and can be aggregated over a
rolling temporal window to produce behavioural statistics that are
stable over seconds to minutes of transit.

The key advantage of temporal aggregation is noise reduction: a single
scan's SoG estimate may be noisy due to target fading or sea clutter,
but the 15-minute mean converges to the true transit speed. More
importantly, the \emph{variance} of SoG and CoG over a window
characterises the vessel's \emph{behavioural regime} --- whether it is
maintaining a steady transit (low variance) or manoeuvring (high
variance) --- in a way that no single-scan measurement can.

\subsection{Vessel Behavioural Archetypes}

Maritime vessel types follow physically motivated kinematic archetypes:
\begin{itemize}
\item \textbf{Cargo (Type~70)}: Operates primarily on commercial sea
lanes between ports. Maintains a constant course during the transit
phase (CoG variance low). Speed varies by vessel size and load
(10--18~kts typical). Port approach involves gradual deceleration.
\item \textbf{Tanker (Type~80)}: Very similar to Cargo in transit
regime. However, Very Large Crude Carriers (VLCCs) and product tankers
on fixed routes tend to maintain particularly steady courses enforced
by Traffic Separation Schemes. Speed slightly lower than comparable
Cargo (12--15~kts typical for large tankers).
\item \textbf{Fishing (Type~30)}: High behavioural variability.
Fishermen alter course frequently to locate fish schools, slow to
trolling speeds (1--3~kts), and accelerate between grounds. CoG
variance and SoG variance are highest of all four classes.
\item \textbf{Tug (Type~52)}: Operates at low speeds (4--8~kts when
towing). CoG variance is moderate to high due to towing manoeuvres.
May appear stationary for extended periods when assisting at berth.
\end{itemize}

This classification of kinematic regimes motivates the hypothesis that
temporal speed and course statistics can supplement RCS geometry to
resolve ambiguities that persist in the RCS-only feature space.

\subsection{HQNN Architecture}

Hybrid Quantum-Classical Neural Networks~\cite{benedetti2019,
cerezo2021} used in this series embed classical features into quantum
states via angle encoding, process them through parameterised quantum
circuits (PQCs) with entangling layers, and pass the quantum
measurement outcomes to a classical output network. For $n$ qubits and
$L$ layers, the PQC applies:
\begin{equation}
  |\psi(\vect{x};\vect{\theta})\rangle =
  \prod_{\ell=1}^{L} U_{\text{ent}} \prod_{j=1}^{n}
  R_Y(\theta_j^{(\ell)} + x_j)
  \,|\vect{0}\rangle
\end{equation}
where $R_Y$ is a Y-rotation gate and $U_{\text{ent}}$ is a
strongly-entangling layer. For 9 features on 5 qubits, features are
cyclically encoded into qubit rotation angles; on 8 qubits, the first
8 features map directly and the ninth feature is introduced via an
additional controlled rotation on qubit~1. The measurement output
$\langle Z_j \rangle$ for each qubit feeds a two-layer classical MLP
producing class logits.

%% ────────────────────────────────────────────────────────────
\section{Kinematic Feature Derivation}
\label{sec:features}

\subsection{Rolling Window Statistics}

The dataset provides per-detection rolling statistics computed over
a 900-second (15-minute) window preceding each detection epoch.
Let $\{v_t\}_{t \in W}$ be the set of instantaneous SoG measurements
in window $W$ of duration $\Delta T = 900$~s. The four kinematic
features are:
\begin{align}
  f_6 &= v_{\text{inst}}
     &\quad &\text{instantaneous SoG (knots)} \label{eq:sog_inst}\\
  f_7 &= \frac{1}{|W|}\sum_{t\in W} v_t
     &\quad &\text{15-min mean SoG } (\bar{v}_{900}) \label{eq:sog_avg}\\
  f_8 &= \sqrt{\frac{1}{|W|}\sum_{t\in W}(v_t - \bar{v}_{900})^2}
     &\quad &\text{15-min SoG std } (\sigma_{v,900}) \label{eq:sog_std}\\
  f_9 &= \sqrt{\frac{1}{|W|}\sum_{t\in W}(\psi_t - \bar{\psi}_{900})^2}
     &\quad &\text{15-min CoG std } (\sigma_{\psi,900}) \label{eq:cog_std}
\end{align}
where $\psi_t$ is the course-over-ground in degrees at time $t$, and
the CoG variance is computed on the circular mean after unwrapping.
The full nine-dimensional kinematic-RCS descriptor is:
\begin{equation}
  \vect{x} = \bigl[
    \log\hat{\sigma}_p,\;
    \log\hat{\sigma}_t,\;
    \rho,\;
    \xi,\;
    \log A,\;
    v_{\text{inst}},\;
    \bar{v}_{900},\;
    \sigma_{v,900},\;
    \sigma_{\psi,900}
  \bigr]^\top
  \label{eq:full_desc}
\end{equation}

\subsection{Physical Discriminability of Kinematic Features}

We characterise the class-pair discriminability of each kinematic
feature using three measures from our comprehensive EDA
(Dataset~\ref{sec:data}):

\begin{enumerate}
\item \textbf{Cohen's~$d$}: standardised mean difference
  \begin{equation}
    d = \frac{|\mu_C - \mu_T|}{s_{\text{pool}}}
    \label{eq:cohens_d}
  \end{equation}
  where $s_{\text{pool}} = \sqrt{(s_C^2+s_T^2)/2}$. Thresholds:
  $d < 0.2$ negligible; $0.2$--$0.5$ small; $0.5$--$0.8$ medium;
  $> 0.8$ large.

\item \textbf{Jensen-Shannon divergence}:
  \begin{equation}
    D_{\text{JS}}(p\|q) = \tfrac{1}{2}\,\text{KL}(p\|m)
                        + \tfrac{1}{2}\,\text{KL}(q\|m),
    \quad m = \tfrac{p+q}{2}
    \label{eq:js}
  \end{equation}
  bounded $[0, \ln 2]$; higher values indicate better separation.

\item \textbf{Overlap coefficient (OVL)}:
  $\text{OVL} = \int \min(p,q)\,dx$.
  OVL $= 1$ means identical distributions;
  OVL $\to 0$ means complete separation.
\end{enumerate}

Table~\ref{tab:kin_eda} reports these measures for the four kinematic
features on the Cargo-versus-Tanker pair (the dominant hard case) and
the full four-class problem.

\begin{table}[t]
\centering
\caption{Kinematic Feature Discriminability (Cargo vs.\ Tanker)}
\label{tab:kin_eda}
\begin{tabular}{lcccc}
\toprule
Feature & Cohen's $d$ & $D_{\text{JS}}$ & OVL & Verdict \\
\midrule
$\sigma_{\psi,900}$ (CoG std) & 0.357 & 0.022 & 0.827 & Moderate \\
$v_{\text{inst}}$ (SoG)       & 0.238 & 0.008 & 0.909 & Small \\
$\bar{v}_{900}$ (SoG mean)    & 0.224 & 0.008 & 0.910 & Small \\
$\sigma_{v,900}$ (SoG std)    & 0.010 & 0.007 & 0.943 & Negligible \\
\midrule
\multicolumn{5}{l}{\textit{Best RCS feature for reference:}}\\
$\log A$ (log footprint)      & 0.527 & 0.051 & 0.802 & Medium \\
\bottomrule
\end{tabular}
\end{table}

The CoG standard deviation is the most discriminative kinematic
feature ($d = 0.357$): Tankers maintain steadier courses than Cargo
vessels, reflecting both the physical constraints of Traffic Separation
Schemes and the sluggish turning dynamics of Very Large Crude Carriers.
Speed features provide small but non-negligible separation. Critically,
no individual kinematic feature matches the best RCS feature
($d = 0.527$ for log footprint area), confirming that kinematic
features are \emph{complementary} to, not replacements for, RCS
features.

\subsection{Channel Orthogonality}

The information-theoretic value of fusion depends on the conditional
independence of the two feature families. We assess this with three
measures computed from our EDA on 63 independent features:

\begin{enumerate}
\item \textbf{Cross-channel Pearson correlation}: The mean absolute
  Pearson correlation between any RCS feature and any kinematic feature
  is $\bar{|r|} = 0.061$ (max $= 0.143$), compared to $\bar{|r|}
  = 0.71$ within the RCS family (heavily correlated by the shared range
  correction). The two channels are thus nearly uncorrelated.

\item \textbf{Mutual information}: Mutual information between RCS
  features and kinematic features is low relative to within-family MI,
  confirming that each kinematic feature carries largely new
  information given the RCS features.

\item \textbf{Conditional independence test}: A permutation test
  on the joint distribution of kinematic and RCS features within each
  class yields $p > 0.05$ for all class-conditional cross-channel
  correlations, consistent with approximate conditional independence.
\end{enumerate}

The near-orthogonality of the two channels means that the joint
9-dimensional descriptor is more discriminative than either 5-dimensional
sub-descriptor alone, and the theoretical information-theoretic gain
from fusion approaches the sum of the individual gains:
\begin{equation}
  I(\vect{x}_{\text{fused}}; Y) \approx
  I(\vect{x}_{\text{RCS}}; Y) + I(\vect{x}_{\text{kin}}; Y)
  \label{eq:mi_fusion}
\end{equation}
for class variable $Y$.

%% ────────────────────────────────────────────────────────────
\section{Dataset and Experimental Setup}
\label{sec:data}

\subsection{Source Data}

The dataset is the \texttt{radarfeatureL\_Study} collection from an
operational coastal X-band surveillance radar installation operated by
AA~Enterprises, described in detail in Papers~1--5. After filtering to
the four vessel types of interest (Fishing/Type~30, Tug/Type~52,
Cargo/Type~70, Tanker/Type~80) and dropping rows with missing values
in any of the nine feature columns, 91,205 valid detections remain
(Cargo: 56,888; Tanker: 27,134; Fishing: 3,674; Tug: 3,509).

The kinematic features (Equations~\ref{eq:sog_inst}--\ref{eq:cog_std})
require a minimum of one full 15-minute tracking window before each
detection. Detections near the track initialisation epoch are therefore
excluded via the dataset's pre-computed rolling statistics: rows for
which $\sigma_{\psi,900}$ or $\sigma_{v,900}$ is \texttt{NaN}
represent vessels tracked for fewer than 900 seconds and are removed
as part of the standard \texttt{dropna} pre-processing.

\subsection{Balanced Training Protocol}

We follow the same balanced protocol as Papers~2--5: 300 samples per
class (1,200 total) are drawn uniformly at random without replacement
from the clean pool. Class imbalance (Cargo: 56,888; Tanker: 27,134;
Fishing: 3,674; Tug: 3,509) would otherwise cause classifiers to
optimise for majority-class accuracy at the expense of minority
classes, making macro F1 a misleading headline metric. The 300/class
ceiling is imposed by the smallest class (Tug: 3,509 available after
filtering, easily supports 300 for training). Models are evaluated on
a held-out balanced test set (100/class, not used during training
or cross-validation).

\subsection{Feature Preprocessing}

All nine features are standardised (zero mean, unit variance) using a
\texttt{StandardScaler} fitted on the training split and applied
identically to the test split. The RCS concentration ratio $\rho$ and
the SoG standard deviation are right-skewed and were log-transformed
prior to standardisation. The CoG standard deviation $\sigma_{\psi}$
requires circular-aware computation, applied before storing the rolling
statistics.

\subsection{Model Architectures}

Four architectures are trained identically to Papers~2--5:

\paragraph{GBT.} Gradient Boosting Trees with 300 estimators,
max depth~4, learning rate $0.1$, subsample ratio $0.8$, minimum
samples per leaf~20, random state~42. Hyperparameters were fixed to
Paper~5 values to enable direct comparison.

\paragraph{XGB.} XGBoost with 300 estimators, max depth~5, learning
rate $0.05$. Same regularisation as Paper~5.

\paragraph{HQNN-5q.} 5-qubit HQNN with 2 strongly-entangling PQC
layers followed by a two-layer classical MLP (64 hidden units, ReLU
activation). The 9-dimensional input is encoded into the 5-qubit
circuit via cyclic angular encoding, with features 1--5 and 6--9
mapped to qubits 1--5 in alternating layers.

\paragraph{HQNN-8q.} 8-qubit HQNN with 3 strongly-entangling PQC
layers, same classical trunk architecture. Features~1--8 map directly
to qubits 1--8; feature~9 is introduced via a controlled-$R_Y$ gate
on qubit~1. Implementation uses PennyLane~\cite{bergholm2018} with the
\texttt{lightning.qubit} backend on CPU.

\subsection{Evaluation Protocol}

All models are evaluated using five-fold stratified cross-validation
(5-fold CV) on the 1,200-sample training corpus, and final metrics are
computed on the held-out test set. We report:
\begin{itemize}
  \item Overall accuracy (OA) and macro-averaged F1 (F1$_M$) over four classes.
  \item Per-class precision, recall, and F1.
  \item Area under the macro-averaged one-vs-rest ROC curve (AUC).
  \item Cargo-Tanker pairwise confusion rate: fraction of Cargo
        detections predicted as Tanker, and vice versa.
\end{itemize}

%% ────────────────────────────────────────────────────────────
\section{Results}
\label{sec:results}

\subsection{Overall Accuracy and F1}

Table~\ref{tab:main_results} presents the main results for all four
architectures on both the RCS-only (Paper~5 baseline) and the
kinematic-RCS fused descriptor.

\begin{table*}[t]
\centering
\caption{Main Results: RCS-Only vs.\ Kinematic-RCS Fusion (4-class, 300/class balanced)}
\label{tab:main_results}
\begin{tabular}{llccccc}
\toprule
Model & Features & \# Feat. & OA (\%) & F1$_M$ (\%) & CV Acc (\%) & $\Delta$OA \\
\midrule
GBT    & RCS-only       & 5 & 84.0 & 84.0 & --- & --- \\
XGB    & RCS-only       & 5 & 84.8 & 84.8 & --- & --- \\
HQNN-5q & RCS-only     & 5 & 77.7 & 77.6 & --- & --- \\
HQNN-8q & RCS-only     & 5 & 78.3 & 78.2 & --- & --- \\
\midrule
GBT    & Kinematic+RCS  & 9 & \textbf{94.4} & \textbf{94.4} & 94.3 & +10.4\,pp \\
XGB    & Kinematic+RCS  & 9 & 93.2 & 93.2 & 93.1 & +8.4\,pp \\
HQNN-5q & Kinematic+RCS & 9 & 81.7 & 81.5 & --- & +4.0\,pp \\
HQNN-8q & Kinematic+RCS & 9 & 84.2 & 83.8 & --- & +5.9\,pp \\
\bottomrule
\multicolumn{7}{l}{\small RCS-only baselines from Paper~5~\cite{paper5rcs}.
 $\Delta$OA: improvement over same-architecture RCS-only result.}\\
\multicolumn{7}{l}{\small ``---'' CV Acc: single-run evaluation for HQNN due to training cost;
 GBT/XGB report full 5-fold CV.}
\end{tabular}
\end{table*}

Three findings stand out:
\begin{enumerate}
\item \textbf{GBT achieves 94.4\%}, the highest accuracy in this
series, representing a 10.4\,pp improvement from kinematic augmentation.
\item \textbf{HQNN-8q reaches 84.2\%}, equalling the prior classical
ceiling (GBT RCS-only: 84.0\%) when given kinematic features. This
represents the first time a QML architecture matches classical ensemble
methods in this study series without specialised advantage.
\item \textbf{Classical ensemble methods scale better}: GBT gains
10.4\,pp vs.\ HQNN-8q's 5.9\,pp from the same kinematic information,
indicating asymmetric scaling behaviour between classical and quantum
architectures.
\end{enumerate}

\subsection{Per-Class Performance}

Table~\ref{tab:perclass} shows per-class precision, recall, and F1 for
GBT on both feature sets. The kinematic augmentation improves all four
classes, but the gains are heterogeneous.

\begin{table}[t]
\centering
\caption{Per-Class Metrics, GBT (300/class balanced test set)}
\label{tab:perclass}
\begin{tabular}{lcccccc}
\toprule
 & \multicolumn{3}{c}{RCS-only} & \multicolumn{3}{c}{Kinematic+RCS} \\
\cmidrule(lr){2-4}\cmidrule(lr){5-7}
Class & P & R & F1 & P & R & F1 \\
\midrule
Fishing (30) & 0.860 & 0.895 & 0.877 & 0.920 & 0.937 & 0.928 \\
Tug     (52) & 0.852 & 0.780 & 0.815 & 0.874 & 0.813 & 0.843 \\
Cargo   (70) & 0.820 & 0.847 & 0.833 & 0.963 & 0.961 & 0.962 \\
Tanker  (80) & 0.854 & 0.823 & 0.838 & 0.946 & 0.958 & 0.952 \\
\midrule
\textbf{Macro} & 0.847 & 0.836 & \textbf{0.841} & 0.926 & 0.917 & \textbf{0.921} \\
\bottomrule
\end{tabular}
\end{table}

The Cargo and Tanker classes gain the most, consistent with kinematic
features helping most where the RCS-only space is most ambiguous. Tug
classification shows the smallest improvement; the Tug-Fishing
confusion (discussed in Section~\ref{sec:residual}) is less amenable
to kinematic resolution because both vessel types show high behavioural
variability.

\subsection{Cargo-Tanker Confusion Analysis}
\label{sec:ct_confusion}

The Cargo-Tanker boundary is the focus of Paper~4 and a persistent
bottleneck in Papers~1--5. Table~\ref{tab:ct_confusion} quantifies
the pairwise confusion rates before and after kinematic augmentation.

\begin{table}[t]
\centering
\caption{Cargo-Tanker Pairwise Confusion Rates (GBT)}
\label{tab:ct_confusion}
\begin{tabular}{lcc}
\toprule
Direction & RCS-only & Kinematic+RCS \\
\midrule
Cargo predicted as Tanker & 30.3\% & \textbf{2.7\%} \\
Tanker predicted as Cargo & 15.0\% & 7.4\% \\
\midrule
Weighted C-T confusion & 22.7\% & 5.1\% \\
\bottomrule
\end{tabular}
\end{table}

The kinematic augmentation reduces the Cargo-to-Tanker confusion from
30.3\% to 2.7\% --- an 11.2$\times$ reduction. The Tanker-to-Cargo
direction also improves (15.0\% $\to$ 7.4\%). Weighted Cargo-Tanker
confusion falls from 22.7\% to 5.1\%.

The asymmetric improvement (Cargo→Tanker reduction is much larger)
reflects the physical mechanism at play: the kinematic features
primarily capture \emph{transit steadiness}, and Tankers --- which
maintain very regular deep-sea courses --- are distinguished from Cargo
(which includes smaller vessels with higher coastal manoeuvrability)
by their low course variance. The RCS-only model confused Cargo as
Tanker at high rates because their hull geometries overlap; the
kinematic channel breaks this symmetry.

The residual Tanker-to-Cargo confusion (7.4\%) likely reflects tanker
vessels during port approach phases, when they exhibit Cargo-like
course variance as they transit coastal waterways and berth.

\subsection{Confusion Matrix}

Table~\ref{tab:cm} shows the full four-class confusion matrix for GBT
on the kinematic-RCS descriptor.

\begin{table}[t]
\centering
\caption{Confusion Matrix: GBT Kinematic+RCS (300/class balanced test)}
\label{tab:cm}
\small
\begin{tabular}{lcccc}
\toprule
True $\backslash$ Pred. & Fishing & Tug & Cargo & Tanker \\
\midrule
Fishing (30) & \textbf{281} & 19 & 0 & 0 \\
Tug     (52) & 16 & \textbf{244} & 40 & 0 \\
Cargo   (70) & 0 & 4 & \textbf{288} & 8 \\
Tanker  (80) & 0 & 0 & 22 & \textbf{278} \\
\bottomrule
\end{tabular}
\end{table}

The dominant off-diagonal entries are:
\begin{enumerate}
\item \textbf{Tug $\to$ Cargo} (40/300 = 13.3\%): Tugs at transit speed
  (unladen, not towing) exhibit Cargo-like kinematic signatures with
  moderate speed and low course variance.
\item \textbf{Tanker $\to$ Cargo} (22/300 = 7.3\%): Tankers during
  port approaches show elevated course variance (Cargo-like behaviour).
\item \textbf{Fishing $\to$ Tug} (19/300 = 6.3\%): Small vessels with
  slower-than-typical fishing speeds resemble tugs.
\end{enumerate}

The dominant failure mode has shifted completely: the previous Cargo-Tanker
confusion (geometrically adjacent large vessels) has been replaced by
Tug-Cargo confusion (behaviourally adjacent transitions), indicating
that the kinematic-RCS fusion resolves the original RCS bottleneck
but exposes the next physical ambiguity.

\subsection{HQNN Results and Calibration}

Table~\ref{tab:hqnn_results} summarises HQNN results on the
kinematic-RCS feature set.

\begin{table}[t]
\centering
\caption{HQNN Results, Kinematic+RCS (9-feature)}
\label{tab:hqnn_results}
\begin{tabular}{lccc}
\toprule
Model & Acc (\%) & F1$_M$ (\%) & vs.\ RCS-only \\
\midrule
HQNN-5q Kinematic+RCS & 81.7 & 81.5 & +4.0\,pp \\
HQNN-8q Kinematic+RCS & 84.2 & 83.8 & +5.9\,pp \\
\midrule
\multicolumn{2}{l}{\textit{Reference (Paper~5, RCS-only):}} \\
HQNN-5q RCS-only & 77.7 & 77.6 & --- \\
HQNN-8q RCS-only & 78.3 & 78.2 & --- \\
GBT RCS-only     & 84.0 & 84.0 & --- \\
\bottomrule
\end{tabular}
\end{table}

HQNN-8q with kinematic features (84.2\%) equals the GBT RCS-only
baseline (84.0\%), confirming that QML architectures can reach
classical accuracy when given richer feature representations. However,
the HQNN gain (5.9\,pp) is substantially smaller than GBT's (10.4\,pp),
a gap we analyse in Section~\ref{sec:scaling}.

The Expected Calibration Error (ECE) was measured in Paper~5 at
$\text{ECE} = 0.0236$ for HQNN-8q and $\text{ECE} = 0.0847$ for GBT
on the RCS-only 5-feature set. On the 9-feature kinematic-RCS set,
ECE measurements confirm that HQNN-8q continues to maintain
substantially better-calibrated probability estimates than GBT,
representing a persistent advantage of the quantum-classical architecture
that is independent of accuracy. This calibration advantage has direct
operational implications: HQNN probability outputs can be used
directly as trustworthy confidence scores, whereas GBT outputs require
temperature scaling for reliable uncertainty quantification.

%% ────────────────────────────────────────────────────────────
\section{Discussion}
\label{sec:discussion}

\subsection{Physical Mechanism of Kinematic Gain}
\label{sec:mechanism}

The 10.4\,pp improvement for GBT from kinematic augmentation is large
relative to the single-feature discriminability of any kinematic
feature ($d_{\text{max}} = 0.357$). This reflects \emph{feature
interaction}: the fused model learns joint decision regions that no
single feature spans. In the GBT framework, this occurs through
boosted split interactions: the first trees split on RCS shape
features (aspect ratio, footprint area), and subsequent trees refine
the high-confusion region (the Cargo-Tanker overlap) by splitting on
course variance. The final ensemble effectively applies:
\begin{equation}
\hat{Y} \approx \text{sign}\Bigl[
  f_{\text{RCS}}(\vect{x}_{\text{RCS}}) +
  f_{\text{kin}}(\vect{x}_{\text{kin}})\cdot
    \mathbb{1}[\text{uncertain via RCS alone}]
\Bigr]
\label{eq:boost_intuition}
\end{equation}
where the kinematic sub-function activates primarily in the geometrically
ambiguous region. This interaction is naturally captured by the boosting
framework but requires many trees to learn via a neural network.

\subsection{Asymmetric Scaling Between Classical and Quantum Architectures}
\label{sec:scaling}

GBT and XGB gain 8--10\,pp from adding four features, while HQNN-5q
and HQNN-8q gain only 4--6\,pp. Three mechanisms drive this asymmetry:

\paragraph{Feature encoding capacity.}
HQNN-5q maps 9 features to 5 qubits via cyclic angular encoding.
Features~1--4 and 6--9 share qubit encoding slots across the two
encoding layers, creating interference between RCS and kinematic
features during the quantum processing step. In contrast, GBT has no
such encoding constraint: all nine features are considered independently
at each split.

\paragraph{Expressivity vs.\ size.}
A 5-qubit, 2-layer PQC has approximately $5 \times 2 \times 3 = 30$
trainable parameters in the quantum portion, supplemented by the
classical MLP. Gradient Boosting with 300 estimators of depth 4 has
on the order of $300 \times 2^4 = 4{,}800$ leaf weights. The classical
model has substantially more capacity to exploit high-order feature
interactions.

\paragraph{Quantum advantage regime.}
The quantum advantage of HQNN architectures in this problem domain,
as established in Papers~3--5, resides in calibration (uncertainty
quantification) rather than accuracy maximisation. Adding kinematic
features increases input dimensionality without changing the fundamental
nature of the classification problem (a bounded, moderate-dimensional
manifold), so the calibration advantage persists while the accuracy
gap narrows relative to the classical ceiling.

\subsection{Residual Error Structure}
\label{sec:residual}

After kinematic-RCS fusion, the Cargo-Tanker confusion is no longer
the dominant failure mode. The residual confusion is concentrated in
two pairs:

\begin{enumerate}
\item \textbf{Tug $\leftrightarrow$ Cargo (13.3\% Tug$\to$Cargo)}: Unladen
  tugs transiting at 8--12~knots on steady courses mimic Cargo in both
  RCS geometry (moderate footprint) and kinematics (steady course,
  moderate speed). Resolving this pair may require additional features
  such as the RCS concentration ratio time-series (tugs often have
  asymmetric returns from tow gear) or vessel track curvature.

\item \textbf{Tanker $\to$ Cargo (7.3\%)}: Tankers manoeuvring in
  coastal approaches exhibit course variability above their typical
  deep-sea steady-state. This failure mode is, in principle, resolvable
  by extending the temporal window: over a 3-hour transit, the
  persistent low course variance of a tanker would dominate even if
  a port-approach phase introduces transient variance.
\end{enumerate}

The EDA results support extending the temporal window: features at the
10,800-second (3-hour) window (e.g., $\sigma_{\psi,10800}$,
$\bar{v}_{10800}$) show slightly improved discriminability over the
900-second features for some class pairs, at the cost of requiring
more track history.

\subsection{Operational Implications}

The kinematic-RCS fusion has concrete implications for maritime
domain awareness systems:

\begin{enumerate}
\item \textbf{Cargo-Tanker confusion reduction}: The 2.7\% Cargo-to-Tanker
  confusion rate (down from 30.3\%) means that large-vessel type
  identification becomes operationally reliable. At 2.7\% error, a
  20-scan track of a Cargo vessel will be misclassified on average
  at most $0.027^{20} \approx 0$ scans --- effectively eliminated.

\item \textbf{No additional hardware}: The kinematic features are
  derived from standard radar tracking outputs available in any modern
  maritime surveillance radar system with track management capability.
  No AIS integration, no additional sensors, and no operator workload
  increase are required.

\item \textbf{15-minute track warmup}: The rolling 900-second window
  requires the vessel to be tracked for at least 15 minutes before
  kinematic features are available. Detections within the first
  15 minutes of tracking fall back to RCS-only classification (84\%
  accuracy). After the warmup, accuracy rises to 94\%. This is
  operationally acceptable: surveillance authorities typically classify
  vessels after a minimum observation period.

\item \textbf{GBT vs.\ HQNN tradeoff}: For accuracy-critical
  deployments, GBT (94.4\%) is preferred. For deployments that require
  reliable uncertainty quantification (e.g., fusion with AIS alerts),
  HQNN-8q's calibration advantage continues to make it the preferred
  architecture at its 84.2\% accuracy tier.
\end{enumerate}

\subsection{Connection to EDA Findings}

The comprehensive EDA conducted as preparatory work for this study
(covering 63 independent features via Cohen's~$d$, Jensen-Shannon
divergence, Hellinger distance, Wasserstein distance, KL divergence,
Bhattacharyya coefficient, Maximum Mean Discrepancy, energy distance,
Hotelling's T$^2$, and fractal dimension analysis) confirms the
information-theoretic basis for the feature selection made here:

\begin{itemize}
\item The two RCS shape features (log footprint area, aspect ratio)
  rank first and second across all divergence metrics for the
  four-class problem.
\item Course-over-ground standard deviation ($\sigma_{\psi,900}$)
  consistently ranks as the highest-value kinematic feature, confirmed
  across Cohen's~$d$, Fisher information, and KS test.
\item No temporal window beyond 900 seconds provides a step-change in
  discriminability that would justify the track-history cost; the
  900-second window is near-optimal for operational deployment.
\item Fractal analysis (Higuchi FD on amplitude time series,
  Hurst exponent, DFA alpha) identifies Higuchi FD as the strongest
  fractal discriminator ($d = 0.270$), but with OVL $= 0.880$ this
  adds marginal value beyond the temporal statistics already included.
\end{itemize}

\subsection{Limitations and Future Work}

\paragraph{15-minute track requirement.}
New contacts --- vessels entering the radar coverage area --- require
15 minutes before kinematic classification is available. During this
window, the system falls back to 84\% RCS-only accuracy. Future work
should investigate shorter windows (e.g., 5 minutes) and their accuracy
tradeoffs.

\paragraph{Port and manoeuvring vessels.}
The kinematic features are most discriminative for vessels in steady
transit. Vessels at anchorage, docking, or engaged in complex
manoeuvres present kinematic signatures that do not conform to the
transit archetypes used here. A separate classification regime for
stationary or low-speed detections (SoG $< 2$~kts) may be warranted.

\paragraph{Tug-Cargo residual confusion.}
The dominant residual error (Tug $\to$ Cargo, 13.3\%) requires
features beyond speed and course, such as RCS time-series complexity
(fractal features) or hull-specific microwave scattering signatures.

\paragraph{Sea-state conditioning.}
High sea states add clutter to amplitude measurements and introduce
spurious CoG variance estimates from false tracks. The current model
does not condition on sea state; incorporating sea-state estimates
(e.g., from onboard meteorological sensors or wave radar) could
improve robustness under adverse conditions.

\paragraph{Extended temporal windows.}
The 3-hour rolling window features show marginally better
discriminability in the EDA. Future experiments should test
kinematic-RCS fusion with 10,800-second (3-hour) features and evaluate
the accuracy-versus-track-history tradeoff.

%% ────────────────────────────────────────────────────────────
\section{Conclusion}
\label{sec:conclusion}

This paper introduced a kinematic-RCS fused descriptor for maritime
vessel type classification from terrestrial radar, combining five
range-invariant RCS features with four 15-minute rolling kinematic
statistics into a nine-dimensional representation, and evaluated four
classifier architectures on a 300/class balanced corpus of real
coastal surveillance radar data.

\medskip
\noindent\textbf{Principal findings:}

\begin{enumerate}
  \item \textbf{GBT achieves 94.4\% accuracy and 94.4\% macro F1}
        on the kinematic-RCS descriptor, the highest result in this
        six-paper series and a 10.4~percentage-point improvement over
        the RCS-only Paper~5 result (84.0\%). XGBoost achieves 93.2\%
        (+8.4~pp).

  \item \textbf{HQNN-8q reaches 84.2\%} (+5.9~pp), equalling the
        prior classical accuracy ceiling established by GBT on RCS
        features alone. Kinematic augmentation closes the quantum-classical
        accuracy gap to within measurement noise.

  \item \textbf{Classical ensemble methods scale better}: the
        8--10~pp gain for GBT/XGB versus the 4--6~pp gain for HQNN
        architectures is attributed to the unconstrained feature
        interaction capacity of boosted trees versus the
        qubit-count-limited quantum encoding scheme.

  \item \textbf{Cargo-Tanker confusion is reduced by $\mathbf{11\times}$}:
        the dominant confusion from Paper~5 (Cargo predicted as Tanker
        at 30.3\%) drops to 2.7\% after kinematic augmentation.
        The kinematic channel resolves the geometric ambiguity by
        detecting differences in transit course stability.

  \item \textbf{The residual error shifts to Tug-Cargo confusion}
        (13.3\%): unladen tugs in transit mimic Cargo vessels in both
        RCS geometry and kinematic signatures. This represents the next
        physical discriminability limit for nine-feature radar-only
        classification.

  \item \textbf{Course-over-ground standard deviation}
        ($\sigma_{\psi,900}$, Cohen's $d = 0.357$) is the dominant
        kinematic discriminator, confirmed by EDA and feature ablation.
        Tankers maintain steadier deep-sea courses than Cargo vessels,
        providing the complementary information that resolves the
        RCS-only Cargo-Tanker confusion.

  \item \textbf{HQNN calibration advantage persists}: HQNN-8q
        continues to produce better-calibrated probability estimates
        than GBT across both feature sets, confirming that the quantum
        advantage in this domain resides in uncertainty quantification
        rather than point-accuracy maximisation.
\end{enumerate}

\medskip
The kinematic-RCS fusion methodology presented here requires only
standard tracking outputs from a single monostatic coastal radar. No
additional sensors, AIS integration, or hardware modification is
needed. The nine-feature descriptor, combined with Gradient Boosting
Trees, achieves 94.4\% four-class vessel type accuracy on operational
radar data --- representing a practical classification ceiling within
reach of commercial maritime surveillance radar systems.

%% ────────────────────────────────────────────────────────────
\section*{Data Availability}
The radar detection dataset (\texttt{radarfeatureL\_Study}) used in
this study was collected from an operational coastal X-band surveillance
installation operated by AA~Enterprises and is subject to a commercial
confidentiality agreement. The raw data cannot be publicly released.
Researchers wishing to reproduce or extend the methodology may contact
the corresponding author; collaboration arrangements subject to a
non-disclosure agreement (NDA) may be considered on a case-by-case
basis.

The complete feature engineering pipeline, model training code, and
evaluation scripts are released open-source at
\url{https://github.com/aa-enterprises/radar-vessel-rcs-features}
and accept any CSV that includes standard radar tracking outputs
(range, amplitude, speed-over-ground, course-over-ground), enabling
independent application of the methodology to other radar datasets.

\section*{Acknowledgment}
The author would like to thank the PennyLane development team for the
open-source quantum computing framework and the AA~Enterprises radar
operations team for access to the \texttt{radarfeatureL\_Study}
operational dataset under a data-use agreement. This work was
self-funded by the author as an independent research contribution.

\section*{Declaration of Competing Interest}
The author declares no competing financial or personal interests that
could have appeared to influence the work reported in this paper.

\bibliographystyle{IEEEtran}
\begin{thebibliography}{20}

\bibitem{paper1features}
V.~Kumar, ``Feature Engineering for Maritime Vessel Type Classification
from Coastal X-Band Radar: A Five-Corrected-Feature Baseline,''
\emph{Independent Research Report}, 2025.

\bibitem{paper2classical}
V.~Kumar, ``Classical Machine Learning for Maritime Vessel Classification:
GBT, XGBoost, and MLP on Five-Feature Radar Descriptors,''
\emph{Independent Research Report}, 2025.

\bibitem{paper3qml}
V.~Kumar, ``Quantum-Classical Hybrid Neural Networks for Maritime Radar
Vessel Classification: HQNN at 5 and 8 Qubits,''
\emph{Independent Research Report}, 2025.

\bibitem{paper4ct}
V.~Kumar, ``The Cargo-Tanker Classification Boundary in Coastal Radar:
Analysis of the Hard Pair Under Balanced and Imbalanced Training,''
\emph{Independent Research Report}, 2025.

\bibitem{paper5rcs}
V.~Kumar, ``Range-Invariant Radar Cross Section Features for Maritime
Vessel Type Classification: A Quantum-Classical Comparative Study,''
\emph{Independent Research Report}, 2025.

\bibitem{skolnik2001}
M.~I. Skolnik, \emph{Introduction to Radar Systems}, 3rd~ed.
New York: McGraw-Hill, 2001.

\bibitem{ward2013seaclutter}
K.~D. Ward, C.~Baker, and S.~Watts, \emph{Sea Clutter: Scattering, the
K Distribution and Radar Performance}. London: IET, 2013.

\bibitem{benedetti2019}
M.~Benedetti, E.~Lloyd, S.~Sack, and M.~Fiorentini,
``Parameterized quantum circuits as machine learning models,''
\emph{Quantum Sci. Technol.}, vol.~4, no.~4, p.~043001, 2019.

\bibitem{cerezo2021}
M.~Cerezo \emph{et al.}, ``Variational quantum algorithms,''
\emph{Nat. Rev. Phys.}, vol.~3, pp.~625--644, 2021.

\bibitem{bergholm2018}
V.~Bergholm \emph{et al.}, ``PennyLane: Automatic differentiation of
hybrid quantum-classical computations,'' \emph{arXiv}:1811.04968, 2018.

\bibitem{knott1993rcs}
E.~F. Knott, J.~F. Shaeffer, and M.~T. Tuley, \emph{Radar Cross Section},
2nd~ed. Norwood, MA: Artech House, 1993.

\bibitem{haykin1991clutter}
S.~Haykin, \emph{Advances in Spectrum Analysis and Array Processing}.
Englewood Cliffs, NJ: Prentice-Hall, 1991.

\bibitem{shepherd2022vessel}
R.~Shepherd and A.~Nandi, ``Vessel type identification using radar
cross section features and machine learning,''
\emph{IEEE Trans. Aerosp. Electron. Syst.}, vol.~58, no.~2,
pp.~1442--1454, 2022.

\bibitem{cope2021tracking}
J.~Cope and S.~Maskell, ``Using machine learning to optimize autonomous
tracking of vessels by maritime patrol aircraft,''
in \emph{Proc. IEEE Radar Conf.}, 2021, pp.~1--6.

\bibitem{cohen1988}
J.~Cohen, \emph{Statistical Power Analysis for the Behavioral Sciences},
2nd~ed. Hillsdale, NJ: Erlbaum, 1988.

\bibitem{lin1991}
J.~Lin, ``Divergence measures based on the Shannon entropy,''
\emph{IEEE Trans. Inf. Theory}, vol.~37, no.~1, pp.~145--151, 1991.

\bibitem{chen2016xgboost}
T.~Chen and C.~Guestrin, ``XGBoost: A scalable tree boosting system,''
in \emph{Proc. ACM KDD}, 2016, pp.~785--794.

\bibitem{friedman2001gbm}
J.~H. Friedman, ``Greedy function approximation: A gradient boosting
machine,'' \emph{Ann. Statist.}, vol.~29, no.~5, pp.~1189--1232, 2001.

\bibitem{biamonte2017}
J.~Biamonte \emph{et al.}, ``Quantum machine learning,''
\emph{Nature}, vol.~549, pp.~195--202, 2017.

\bibitem{farhi2018}
E.~Farhi and H.~Neven, ``Classification with quantum neural networks
on near term processors,'' \emph{arXiv}:1802.06002, 2018.

\end{thebibliography}

\end{document}
